\documentclass[11pt]{article}

\usepackage[margin=1in]{geometry}
\usepackage{amsmath,amssymb,amsthm,mathtools,mathrsfs,booktabs,tabularx,graphicx}
\usepackage{microtype,xurl}
\usepackage{xcolor}
\usepackage[colorlinks=true,linkcolor=blue!55!black,
            citecolor=blue!55!black,urlcolor=blue!55!black]{hyperref}

\newtheorem{theorem}{Theorem}[section]
\newtheorem{proposition}[theorem]{Proposition}
\newtheorem{corollary}[theorem]{Corollary}
\newtheorem{lemma}[theorem]{Lemma}
\theoremstyle{definition}
\newtheorem{definition}[theorem]{Definition}
\theoremstyle{remark}
\newtheorem{remark}[theorem]{Remark}

\newcommand{\R}{\mathbb R}
\newcommand{\C}{\mathbb C}
\newcommand{\dd}{\mathrm d}
\newcommand{\Ric}{\operatorname{Ric}}
\newcommand{\coker}{\operatorname{coker}}
\newcommand{\rank}{\operatorname{rank}}
\newcommand{\inertia}{\operatorname{inertia}}
\newcommand{\cert}[1]{\nolinkurl{#1}}
\newcolumntype{Y}{>{\raggedright\arraybackslash}X}

\title{Lorentzian Endpoint Conditions Do Not Select\\
the Einstein Sector of Pure Weyl Gravity}
\author{GPT-5.6.sol\thanks{OpenAI model and principal programme author.}
\and Claude Fable 5\thanks{Anthropic model and computational coauthor. The
project was commissioned and coordinated by Asger Alstrup Palm
(\texttt{asger@area9.dk}), who initiated the questions, exercised editorial
control, and serves as corresponding human contact, but claims no technical
authorship.}}
\date{Focused working draft, 24 July 2026}

\begin{document}
\maketitle

\begin{abstract}
Maldacena-type boundary selection isolates an Einstein sector of conformal
gravity in Euclidean anti-de Sitter or late-time de Sitter settings.  We ask
whether physically familiar one-ended Lorentzian conditions do the same for
linear perturbations of a Schwarzschild black hole in four-dimensional pure
Weyl gravity.  They do not.

In the axial \(\ell=2\) sector we derive the linearized Bach system from the
Weyl action and identify an exact filtered differential module with successive
quotients
\[
M_{\rm RW}^{(2)},\qquad M_{\rm RW}^{(2)},\qquad M_{\rm RW}^{(1)}.
\]
The two spin-two factors form a non-split self-extension.  The action-derived
Lee--Wald form on the incoming trace space is nondegenerate with signature
\((1,2)\).  In factor coordinates it is congruent, for every real
\(\omega>0\), to
\[
\pi\alpha_{\rm W}
\begin{pmatrix}
0&576\omega/5&0\\
576\omega/5&0&0\\
0&0&-32/(15\omega)
\end{pmatrix}.
\]
Scalar Regge--Wheeler Wronskians and boundary d\'evissage prove that the map
from future-horizon-regular data to this complete incoming Krein space is
invertible for every positive real frequency.  Negative-flux directions are
therefore populated by regular exterior solutions, rather than being formal
endpoint artefacts.

On the certified cell
\(\omega\in[0.49995,0.50005]\) the outgoing and future-horizon spaces are
matching three-channel Krein spaces, the outgoing map is invertible, and the
one-sided scattering map obeys the exact integrated pseudo-isometry.  Analytic
continuation gives outgoing population on an open dense, full-measure subset
of the positive axis; possible failures are isolated scalar reflection zeros.
The same filtration excludes exponentially growing separated axial modes.

Finally, a validated projective Evans calculation encloses one simple damped
spin-two quasinormal frequency and certifies a nonzero intrinsic extension
selector.  The spin-one factor is a local unit there.  The complete connection
therefore has Smith valuations \((0,0,2)\), algebraic multiplicity two,
geometric multiplicity one, and one length-two connection root chain.  This is
a connection-level intrinsic exceptional point.  A second-order physical
Green-resolvent pole remains conditional on a separate analytic Fredholm
realization.
\end{abstract}

\section{Question, answer, and scope}

Pure Weyl gravity is defined by
\begin{equation}
 S_{\rm W}[g]
 =\alpha_{\rm W}\int_M\sqrt{-g}\,
 C_{abcd}C^{abcd}\,\dd^4x,
 \qquad
 B_{ab}=0.
\end{equation}
Every Einstein metric is Bach-flat, but the fourth-order Bach equation has
additional local solutions.  The question addressed here is deliberately
narrow:

\begin{quote}
Do future-horizon regularity and the ordinary incoming or outgoing
real-frequency conditions at null infinity force a Schwarzschild
perturbation of pure Weyl gravity into the linearized Einstein kernel?
\end{quote}

For the complete axial \(\ell=2\) system the answer is no.  The extra
curvature carrier survives the endpoint conditions, participates in the
action-derived current, and supplies populated negative directions in a
nondegenerate Krein space.

This is not a statement that no boundary condition can select the Einstein
sector.  The local Cauchy restriction
\[
\delta R_{ab}|_\Sigma=0,
\qquad
\nabla_n\delta R_{ab}|_\Sigma=0
\]
does select the linearized Einstein kernel in the axial carrier problem.
Nor is this paper a quantum-positivity theorem.  ``Negative'' below always
means negative with respect to the classical Lee--Wald flux form.

\subsection{Five principal results}

The paper is organized around five statements.

\begin{enumerate}
\item The complete axial system has an exact
RW/RW/spin-one filtration, with a nontrivial repeated-spin-two extension.
\item Its incoming action-flux space has signature \((1,2)\) and is populated
at every \(\omega>0\) by future-horizon-regular solutions.
\item On a certified nonzero frequency cell the complete one-sided scattering
map is pseudo-isometric between matching Krein spaces; outgoing population is
generic on the positive real axis.
\item The full filtered system has no exponentially growing separated axial
mode.
\item One enclosed damped spin-two quasinormal frequency is a certified
connection-level exceptional point with Smith valuations \((0,0,2)\).
\end{enumerate}

\subsection{Controlled vocabulary}

\begin{definition}
The \emph{Einstein kernel} is the kernel of the linearized Ricci map
\(h\mapsto\delta R_{ab}[h]\).  The \emph{Ricci carrier} is its image on
linearized Bach solutions.  \emph{Population} of an endpoint trace space
means surjectivity of the relevant horizon-to-endpoint connection map.
\emph{Factor frame} means a basis adapted to the differential-module
filtration.  A \emph{signature frame} diagonalizes a Hermitian flux form.
These two frames are not identified.
\end{definition}

The result tags used below have the repository meanings
\textsc{local-algebraic}, \textsc{reduced-mode}, and
\textsc{lorentzian-causal}.  A reduced-mode theorem is never used as evidence
for a complete Lorentzian quantum or time-domain construction.

\section{Relation to existing work}

Maldacena's conformal-gravity construction uses a Neumann condition in
Euclidean AdS or at late times in de Sitter space to select Einstein
solutions \cite{Maldacena}.  The present problem has a different signature,
background, and boundary geometry: a one-ended Lorentzian Schwarzschild
exterior with a future event horizon and null infinity.

Ricci-tensor variables have long been useful for exposing higher-derivative
spin-two sectors.  Recent quadratic-gravity work computes axial quasinormal
modes of massless and massive spin-two fields \cite{Antoniou2025}, develops a
Ricci-tensor treatment of Kerr perturbations \cite{Knoska2025}, evolves
Einstein--Weyl perturbations in time \cite{Konoplya2025}, and studies the
excitation and suppression of beyond-GR ringdown by point-particle sources
\cite{Antoniou2026}.  A recent thesis gives a complementary radial
Ricci-tensor reconstruction and its relation to the Schwarzschild--Bach
family \cite{Tajcovsky2026}.

Those works concern quadratic or Einstein--Weyl theories, massive sectors,
radial reconstructions, or observable excitation.  Our subject is the
degenerate massless self-extension in strict pure \(C^2\) gravity and the
indefinite current derived from that action.  The claims are therefore
complementary: we do not infer excitation amplitudes or a complete waveform
from endpoint population alone.

\section{Ricci--Bach composition}

Let the background be Ricci-flat.  For a metric perturbation \(h_{ab}\), put
\[
\psi_{ab}=\delta R_{ab}[h],
\qquad
\psi=g^{ab}\psi_{ab}.
\]
Linearization of the Bach tensor gives
\begin{equation}
\label{eq:ricci-bach}
\delta B_{ab}[h]
=
\frac12\Box\psi_{ab}
+C_{acbd}\psi^{cd}
-\frac16\nabla_a\nabla_b\psi
-\frac1{12}g_{ab}\Box\psi.
\end{equation}
In the axial sector \(\psi=0\), so the carrier equation is a second-order
Lichnerowicz-type equation
\begin{equation}
\label{eq:axial-carrier}
L_{\rm car}\psi_{ab}
:=
\frac12\Box\psi_{ab}+C_{acbd}\psi^{cd}=0.
\end{equation}

\begin{proposition}[Ricci--Bach exact sequence]
\label{prop:exact-sequence}
On the declared axial Schwarzschild \(\ell=2\) mode class, the linearized
Ricci map fits into the exact sequence
\begin{equation}
\label{eq:exact-sequence}
0\longrightarrow
\mathcal E_{\rm RW}
\longrightarrow
\mathcal B_{\rm axial}
\overset{\delta\Ric}{\longrightarrow}
\mathcal C_{\rm axial}
\longrightarrow0,
\end{equation}
where \(\mathcal E_{\rm RW}\) is the two-dimensional Regge--Wheeler
Einstein module and \(\mathcal C_{\rm axial}\) is the realized
second-order Ricci carrier.  The sequence is exact but has no certified
canonical metric splitting.
\end{proposition}

\begin{proof}[Proof sketch]
Equation \eqref{eq:ricci-bach} proves that the Ricci image of a Bach solution
obeys \eqref{eq:axial-carrier}.  The complete reconstruction audit verifies
every Ricci row, including the independent \(v\varphi\) equation omitted in an
earlier shallow reduction.  Its zero fibre is exactly the linearized Einstein
Regge--Wheeler system.  Exact recurrence and reconstruction certificates show
that the declared carrier is realized.  A rational splitting is separately
obstructed by the extension invariants discussed in Section~\ref{sec:jet}.
\end{proof}

\subsection{Horizon reach}

In ingoing Eddington--Finkelstein coordinates the carrier horizon is a regular
singular point.  For every nonzero real frequency its ingoing analytic
solution space is two-dimensional.  Thus horizon regularity alone does not
imply \(\psi_{ab}=0\).

\begin{corollary}[Local horizon nonselection]
Future-horizon analyticity admits nonzero Ricci-carrier data in addition to
the Einstein kernel.  Hence it does not select the Einstein image within the
axial Bach solution space.
\end{corollary}

\section{The axial filtration and its partial jet}
\label{sec:jet}

After exact factor reduction the six-state axial system has a filtration
\[
0\subset F_1\subset F_2\subset F_3=M_{\rm axial},
\]
with
\begin{equation}
\label{eq:filtration}
F_1\simeq M_{\rm RW}^{(2)},\qquad
F_2/F_1\simeq M_{\rm RW}^{(2)},\qquad
F_3/F_2\simeq M_{\rm RW}^{(1)}.
\end{equation}
Here \(M_{\rm RW}^{(2)}\) is the axial spin-two Regge--Wheeler module and
\(M_{\rm RW}^{(1)}\) is the spin-one quotient with potential
\[
V_1(r)=\frac{6(r-2)}{r^3}
\]
in units \(M=1\).

In factor variables \(X,Y,Z\in\C^2\), the generator is
\begin{equation}
\label{eq:six-state}
\begin{aligned}
X'&=AX+EY+CZ,\\
Y'&=AY+DZ,\\
Z'&=A_1Z.
\end{aligned}
\end{equation}
The two upper sources share one image line:
\[
E=USJ,\qquad C=USN.
\]
Consequently \eqref{eq:six-state} is exactly the spin-two-row first jet of
the four-state family
\begin{equation}
\label{eq:four-state}
\mathcal B(\tau)=
\begin{pmatrix}
A+\tau E&D+\tau C\\
0&A_1
\end{pmatrix}.
\end{equation}
Over the dual numbers \(\C[\varepsilon]/(\varepsilon^2)\), the six-state
solution is simply \((Y+\varepsilon X,Z)\).

\begin{theorem}[Exact axial partial-jet filtration]
\label{thm:jet}
Every interior transport matrix has the form
\begin{equation}
\label{eq:jet-transport}
\Phi_6=
\begin{pmatrix}
P&\dot P&\dot Q\\
0&P&Q\\
0&0&R
\end{pmatrix}_{\tau=0},
\end{equation}
where
\[
\begin{pmatrix}P(\tau)&Q(\tau)\\0&R\end{pmatrix}
\]
is the transport of \eqref{eq:four-state}.  In particular
\[
\det\Phi_6=(\det P)^2\det R.
\]
The repeated-spin-two extension is non-split over the generic rational
differential field and is not generated, modulo rational gauge, by frequency,
Schwarzschild-mass, or angular-eigenvalue differentiation.
\end{theorem}

The jet identity is useful both conceptually and numerically.  A triangular
connection
\[
T=
\begin{pmatrix}
a&b&c\\0&a&d\\0&0&f
\end{pmatrix}
\]
is the jet of
\[
C(\tau)=\begin{pmatrix}a(\tau)&d(\tau)\\0&f\end{pmatrix},
\quad
\dot a=b,\quad \dot d=c.
\]
Therefore
\begin{equation}
\label{eq:inverse-ratios}
T^{-1}=
\begin{pmatrix}
a^{-1}&-b/a^2&(bd-ac)/(a^2f)\\
0&a^{-1}&-d/(af)\\
0&0&f^{-1}
\end{pmatrix},
\end{equation}
and the three apparently independent extension ratios are one base transfer
and its intrinsic derivative.

\subsection{A local branch-resolving \(C\) operator is obstructed}

The non-split extension also constrains possible positive-metric
constructions.  The statement is local and classical; it is not a theorem
about a quantum Hilbert space.

\begin{proposition}[No branch-resolving rational involution]
\label{prop:no-local-c}
Let
\[
0\longrightarrow M_{\rm RW}
\longrightarrow E_{\rm Bach}
\longrightarrow M_{\rm RW}
\longrightarrow0
\]
be the generic repeated-spin-two extension.  There is no rational
differential-module endomorphism \(C\) with \(C^2=1\) that acts by \(+1\) on
the Einstein submodule and by \(-1\) on the quotient.
\end{proposition}

\begin{proof}
In a triangular frame,
\[
\mathcal A=\begin{pmatrix}A&E\\0&A\end{pmatrix}.
\]
A filtration-compatible involution with the declared actions must be
\[
C=\begin{pmatrix}I&Q\\0&-I\end{pmatrix}.
\]
Covariant constancy,
\[
C'+C\mathcal A-\mathcal AC=0,
\]
requires
\[
Q'+QA-AQ=-2E.
\]
Over characteristic zero this says that the extension cocycle is a rational
coboundary.  It would split the extension, contradicting the certified
nonzero rational extension class.
\end{proof}

There is an independent endpoint obstruction to a branch-preserving positive
fundamental symmetry.  The Einstein spin-two vector \(E\) is null.  If a
fundamental symmetry preserved its line, \(CE=\pm E\), then the proposed
positive majorant would satisfy
\[
G(E,CE)=\pm G(E,E)=0.
\]
Thus any positive fundamental symmetry must mix the two spin-two null
directions.  The elementary endpoint symmetry that exchanges them is a
perfectly valid classical Krein fundamental symmetry, but it is not
Mannheim's dynamical \(C\) operator.  Proposition~\ref{prop:no-local-c}
shows that this distinction is structural: a dynamical \(C\), if it exists,
cannot be a local rational operator that resolves the two layers with
opposite signs.

\section{Action-derived endpoint flux}

The sphere-integrated Lee--Wald current is fixed by the off-shell Green
identity
\begin{equation}
\label{eq:green}
\partial_tF^t(h_A,h_B)+\partial_rF^r(h_A,h_B)
=4\alpha_{\rm W}\int_{S^2}\sqrt q\,
\bigl(
h_B^{ab}\delta B_{ab}[h_A]
-h_A^{ab}\delta B_{ab}[h_B]
\bigr).
\end{equation}
It is therefore not an arbitrarily normalized Wronskian of a chosen master
field.

Let \(G_-(\omega)\) be the incoming null-infinity Gram in the three-dimensional
factor trace space.  Define
\[
E=EI0,\qquad
R=XI0-\frac{i}{\omega}XI1,
\]
and let \(S\) be the unique \(G_-\)-orthogonal lift of the unit spin-one
quotient.  After the harmless Einstein shear making the second spin-two
factor null, the exact factor Gram is
\begin{equation}
\label{eq:factor-gram}
\widehat G_-^{\rm fac}
:=\frac{G_-}{\pi\alpha_{\rm W}}
=
\begin{pmatrix}
0&\dfrac{576}{5}\omega&0\\[2mm]
\dfrac{576}{5}\omega&0&0\\[2mm]
0&0&-\dfrac{32}{15\omega}
\end{pmatrix}.
\end{equation}

\begin{theorem}[Canonical incoming Krein anatomy]
\label{thm:krein}
For every real \(\omega>0\),
\[
\inertia G_-(\omega)=(1,2,0).
\]
The repeated spin-two extension is a hyperbolic plane and the intrinsic
spin-one quotient metric is strictly negative:
\[
g_1=-\pi\alpha_{\rm W}\frac{32}{15\omega}.
\]
In normalized channels
\[
P_2=\frac{E+R}{\sqrt{2a}},\qquad
N_2=\frac{E-R}{\sqrt{2a}},\qquad
N_1=\frac{S}{\sqrt b},
\]
where \(a=576\omega/5\) and \(b=32/(15\omega)\), the Gram is
\[
\operatorname{diag}(1,-1,-1).
\]
\end{theorem}

The positive Hilbert majorant associated with this Krein form is
\begin{equation}
\label{eq:majorant}
\|(e,r,s)\|_{\rm maj}^2
=
\int_0^\infty
\left[
\frac{576}{5}\omega(|e|^2+|r|^2)
+\frac{32}{15\omega}|s|^2
\right]\dd\omega.
\end{equation}
Thus the natural time-domain exponents are
\[
(e,r)\in\dot H^{1/2},
\qquad
s\in\dot H^{-1/2}.
\]
This identifies the canonical completion but does not itself prove global
boundedness of the scattering multiplier.

\section{Every incoming channel is populated}

Let \(T_-(\omega)\) map future-horizon-regular coefficients to incoming
null-infinity factor traces.  Its factor form is triangular:
\begin{equation}
\label{eq:tminus}
T_-=
\begin{pmatrix}
a_2&b&c\\
0&a_2&d\\
0&0&a_1
\end{pmatrix},
\qquad
a_s=A_{{\rm in},s}.
\end{equation}

\begin{lemma}[Boundary d\'evissage]
\label{lem:devissage}
Suppose a filtered differential system has boundary-compatible inclusion and
quotient maps, and none of its scalar quotient problems admits a nonzero
solution in the declared boundary class.  Then the full filtered system has
no nonzero solution in that class.
\end{lemma}

\begin{proof}
Project a full solution to the top quotient.  The quotient boundary problem
forces the projection to vanish, so the solution lies in the preceding
submodule.  Repeat through the filtration.  The last scalar factor then
forces the remaining solution to vanish.  No diagonalization or splitting of
the extension is required.
\end{proof}

For either scalar Regge--Wheeler factor, real-frequency current conservation
gives
\begin{equation}
\label{eq:wronskian}
|A_{{\rm in},s}|^2-|A_{{\rm out},s}|^2=1.
\end{equation}
Thus \(A_{{\rm in},s}\neq0\) for every real \(\omega>0\).

\begin{theorem}[Global incoming population]
\label{thm:incoming}
For the complete axial \(\ell=2\) system,
\[
T_-(\omega)\in GL(3,\C)
\qquad\text{for every real }\omega>0.
\]
Consequently every positive, negative, and null vector of the incoming
Krein space is the trace of a unique future-horizon-regular solution.
\end{theorem}

\begin{proof}
A vector in \(\ker T_-\) would be horizon regular and purely outgoing at
infinity.  Its spin-one projection vanishes by
\eqref{eq:wronskian}.  Boundary d\'evissage then removes the carrier and
metric spin-two factors successively.  Hence the kernel is zero.  Since
\(T_-\) is a \(3\times3\) map, it is invertible.
\end{proof}

The conclusion is stronger than existence of formal negative directions:
the entire indefinite incoming space is physically reached within the
declared exterior boundary problem.

\section{Outgoing population and one-sided Krein scattering}

Let \(T_+\) be the future-null-infinity outgoing connection and
\(H_{\mathcal H^+}\) the outward-oriented future-horizon Gram.  Conservation
of \eqref{eq:green} gives
\begin{equation}
\label{eq:stokes}
H_{\mathcal H^+}
+T_+^\dagger G_+T_+
-T_-^\dagger G_-T_-=0.
\end{equation}
The explicit entries of \(T_+\) are not needed to decide population on the
certified cell.  The scalar outgoing factors and the exact filtration give
\[
\det T_+=\rho_+A_{{\rm out},2}^2A_{{\rm out},1},
\qquad \rho_+\neq0.
\]

\begin{theorem}[Certified outgoing cell]
\label{thm:outgoing-cell}
For \(M=1\), axial \(\ell=2\), and
\[
I_0=[0.49995,0.50005],
\]
validated endpoint calculations prove
\[
|A_{{\rm out},2}|>0.011326240435330133,
\qquad
|A_{{\rm out},1}|>0.14680548488890086.
\]
Hence \(T_+(\omega)\) is invertible throughout \(I_0\).  On the exact common
cell
\[
I_*=[1/2,10001/20000],
\]
the incoming, outgoing, and future-horizon coefficient spaces are
nondegenerate three-channel Krein spaces of inertia \((1,2,0)\).
\end{theorem}

In incoming coordinates define
\[
R=T_+T_-^{-1},
\qquad
A=T_-^{-1}.
\]
Multiplying \eqref{eq:stokes} by \(T_-^{-\dagger}\) and \(T_-^{-1}\) gives
\begin{equation}
\label{eq:pseudo}
G_-
=R^\dagger G_+R+A^\dagger H_{\mathcal H^+}A.
\end{equation}
Equivalently,
\[
\mathscr S=
\begin{pmatrix}R\\A\end{pmatrix}
\]
is a pseudo-isometric embedding
\[
(\C^3,G_-)
\hookrightarrow
(\C^3\oplus\C^3,G_+\oplus H_{\mathcal H^+}).
\]

\begin{corollary}[Band-limited pseudo-isometry]
Multiplication by \(T_+\) is a bounded isomorphism on
\(L^2(I_*;\C^3)\).  Equation \eqref{eq:pseudo} therefore integrates to an
exact pseudo-isometry of the complete band-limited wave-packet spaces.
\end{corollary}

\begin{corollary}[Generic positive-real outgoing population]
\label{cor:generic-outgoing}
Assume the factor-normalized scalar Jost coefficients are holomorphic away
from the separate threshold \(\omega=0\).  Their strict nonvanishing on
\(I_0\) proves that neither is identically zero.  Therefore the positive-real
exceptional set is locally finite and
\[
T_+(\omega)\in GL(3,\C)
\]
on an open dense, full-measure subset of \((0,\infty)\).
On every compact positive-frequency band the corresponding multiplication
operator is injective with dense range.  It is a bounded isomorphism exactly
when the band contains no exceptional frequency.
\end{corollary}

This does not prove that the exceptional set is empty.  The next scattering
experiment is to compute the typed matrix \(T_+\), the reflection map \(R\),
and their extension shears on a substantial real-frequency band using the
partial-jet system \eqref{eq:four-state}.

\section{Exact threshold structure}
\label{sec:threshold}

The threshold \(\omega=0\) is excluded from the preceding holomorphic
genericity argument and must be treated separately.  Several exact facts can
nevertheless be established before a uniform low-frequency scattering
estimate.

For
\[
D=\frac{r-2}{r}\partial_r,\qquad
V_2=\frac{6(r-2)(r-1)}{r^4},\qquad
V_1=\frac{6(r-2)}{r^3},
\]
consider \((D^2-V_s)\phi=0\).

\begin{proposition}[Exact scalar threshold nonresonance]
\label{prop:threshold}
The horizon-normalized regular zero modes are
\[
\phi_{0,2}(r)=\frac{r^3}{8},\qquad
\phi_{0,1}(r)=\frac{r^2(2r-3)}4.
\]
Independent solutions are, up to nonzero constant multiples,
\[
\widetilde\phi_{0,2}
=
\frac{r^3}{32}\log\frac{r-2}{r}
+\frac{r^2}{16}+\frac r{16}+\frac1{12}+\frac1{8r},
\]
and
\[
\widetilde\phi_{0,1}
=
\frac{
3r^2(2r-3)\log\frac{r-2}{r}
+12r^2-6r-2}{24}.
\]
The regular modes grow as nonzero multiples of \(r^3\).  The decaying
companions behave respectively as
\(-1/(5r^2)+O(r^{-3})\) and
\(-1/(10r^2)+O(r^{-3})\), but are logarithmically singular at the horizon.
Consequently neither scalar factor has a bounded zero-energy resonance.
\end{proposition}

\begin{proof}
Direct substitution gives
\[
(D^2-V_s)\phi_{0,s}=0,\qquad \phi_{0,s}(2)=1,
\]
and verifies the two reduction-of-order companions.  Horizon regularity
selects the polynomial solution in each scalar problem; neither selected
solution is bounded at infinity.  Conversely, the solutions that decay at
infinity contain \(\log(r-2)\).  This exhausts each two-dimensional scalar
solution space.
\end{proof}

The exact reduced projective cocycle also exposes the relation between the
three factors:
\begin{equation}
\label{eq:threshold-cocycle}
\mathcal I_{\rm red}
=
\frac{i(r-2)(2r\omega^2+3\omega^2+12)}{5r^4\omega}
=
\frac{2i}{5\omega}(V_1-V_2)
+\frac{i\omega(r-2)(2r+3)}{5r^4}.
\end{equation}
Thus its singular threshold term is exactly the difference of the spin-one
and spin-two potentials.  On \(\phi_{0,2}\) the leading source is
\[
\frac{3i}{10\omega}\frac{r-2}{r}.
\]
It has the elementary spin-two primitive
\begin{equation}
\label{eq:threshold-primitive}
(D^2-V_2)p=\frac{r-2}{r},\qquad
p=-\frac{r^2}{4}-\frac r4-\frac13-\frac1{2r}.
\end{equation}
The horizon-fixed choice
\[
p_H=p+\frac{25}{12}\phi_{0,2}
\]
satisfies \(p_H(2)=0\).

\begin{remark}[Open matching estimate]
Formal two-region matching to the far \(\ell=2\) Bessel equation predicts
\[
|A_{{\rm in},2}|\sim |A_{{\rm out},2}|
\sim\frac{15}{16\omega^3},\qquad
|A_{{\rm in},1}|\sim |A_{{\rm out},1}|
\sim\frac{15}{4\omega^3},
\]
and therefore
\[
|\det T_+|\sim\frac{3375}{1024}\omega^{-9}.
\]
The primitive \eqref{eq:threshold-primitive} further suggests
\(b/a^2=O(\omega^2)\) in the compatible intrinsic normalization.  These
asymptotics are not used as theorems here.  A two-region Volterra remainder
controlling the Schwarzschild tail and an endpoint-compatible extension
normalization are still required before one may infer a punctured
zero-frequency interval of outgoing invertibility or a weighted multiplier
bound.
\end{remark}

\section{Separated-mode stability}

In the Fourier convention used here, exponentially growing modes lie in the
lower half \(\omega\)-plane.  For the scalar spin-one and spin-two
Regge--Wheeler equations, an ingoing-horizon/outgoing-infinity solution in
that half-plane decays at both ends.  Their real, nonnegative potentials give
the standard energy identity, forcing the scalar solution to vanish.

\begin{theorem}[No growing axial separated mode]
\label{thm:no-growth}
The complete axial \(\ell=2\) six-state system has no nonzero separated
solution satisfying the future-horizon-regular and pure-outgoing-infinity
conditions in the growth half-plane.
\end{theorem}

\begin{proof}
The boundary maps preserve the declared complex-frequency classes, including
regular replacements at the apparent frame events.  Apply
Lemma~\ref{lem:devissage}: the spin-one projection vanishes by the scalar
energy identity, then the carrier spin-two quotient vanishes, and finally the
metric spin-two factor vanishes.
\end{proof}

This is separated-mode stability.  It does not supply limiting absorption,
uniform resolvent bounds, threshold control, decay, or bounded time evolution.
The extension ratios in \eqref{eq:inverse-ratios} can still generate
non-normal amplification.

\section{A certified intrinsic connection exceptional point}
\label{sec:ep2}

At a simple spin-two quasinormal frequency \(\omega_n\), assume the spin-one
factor is a local unit.  Analytic elimination of that unit reduces the
repeated block to
\begin{equation}
\label{eq:local-block}
T_2(\omega)=
\begin{pmatrix}
a(\omega)&b(\omega)\\
0&a(\omega)
\end{pmatrix},
\qquad
a(\omega_n)=0,\quad a'(\omega_n)\neq0.
\end{equation}

\begin{proposition}[Local Smith dichotomy]
\label{prop:smith}
If \(b(\omega_n)\neq0\), the full \(3\times3\) connection has Smith
valuations
\[
(0,0,2),
\]
rank two at the QNM, geometric multiplicity one, and one length-two root
chain.  If \(b(\omega_n)=0\), the valuations are
\[
(0,1,1),
\]
the rank is one, and there are two independent root vectors.
\end{proposition}

\begin{proof}
Work in the local discrete valuation ring with
\(z=\omega-\omega_n\).  In the first case \(b\) is a unit, so the ideal of
\(2\times2\) minors is the unit ideal while
\(\det T\) has valuation two.  The invariant factors are therefore
\((0,0,2)\).  In the second case \(b=ag\), every \(2\times2\) minor is
divisible by \(a\), while the minor \(af\) has valuation one; hence
\((0,1,1)\).
\end{proof}

The selector can be computed projectively.  Transport the horizon-regular
and infinity-outgoing scalar lines to a common matching section and use a
Riccati chart \(v=Y_2/Y_1\).  With
\[
\delta(\omega,\tau)=v_\infty(\omega,\tau)-v_H(\omega,\tau),
\]
the scalar Evans coefficient has the form \(a=g\delta\), where \(g\) is an
analytic unit.  At a simple QNM,
\begin{equation}
\label{eq:velocity}
\frac{b(\omega_n)}{a'(\omega_n)}
=
\frac{\delta_\tau(\omega_n)}{\delta_\omega(\omega_n)}
=-\omega_n'(0).
\end{equation}

\begin{theorem}[Certified axial connection EP2]
\label{thm:ep2}
There is a radius-\(10^{-7}\) disk centered at
\[
\omega_c
=-0.3736716844180418
+0.0889623156889357\,i
\]
that contains exactly one simple spin-two quasinormal zero in the declared
\(e^{+i\omega t}\) convention.  On that disk:
\begin{enumerate}
\item the validated projective Evans contour has winding \(+1\);
\item \(\delta_\omega\) excludes zero;
\item the intrinsic tangent satisfies
\[
\delta_\tau\in
[0.0010932\pm7.60\times10^{-8}]
+i[-0.0002195\pm8.89\times10^{-8}],
\]
and therefore excludes zero;
\item the spin-one projective mismatch has modulus greater than \(1/2500\).
\end{enumerate}
Consequently the complete axial connection has Smith valuations
\((0,0,2)\), algebraic multiplicity two, geometric multiplicity one, and one
length-two connection root chain at the unique enclosed QNM.
\end{theorem}

The exceptional point is intrinsic: the rational self-extension is not a
frequency, mass, or angular-eigenvalue tangent modulo admissible rational
gauge.  What is certified is a pole of order two of the inverse
\emph{connection} block.  To infer the rank-one principal part
\[
-\frac{\beta_n}{\alpha_n^2}
\frac{u_n\otimes\widetilde u_n}{(\omega-\omega_n)^2}
\]
for the physical differential resolvent, one must still construct an analytic
Fredholm realization with the QNM endpoint conditions and prove graph-norm
boundedness of the extension.

\section{Polar endpoint reach}

The polar carrier includes trace terms:
\begin{equation}
\delta B_{ab}
=\frac12\Box\psi_{ab}+C_{acbd}\psi^{cd}
-\frac16\nabla_a\nabla_b\psi
-\frac1{12}g_{ab}\Box\psi.
\end{equation}
On the Bianchi-constrained polar \(\ell=2\) system the operator is traceless
and has a one-function linearized conformal-gauge direction.  On a traceless
slice the horizon residue analysis yields, after quotienting the regular
conformal direction, a two-parameter physical ingoing-analytic carrier
family.

\begin{corollary}[Parity-complete local endpoint nonselection]
At the linear \(\ell=2\) mode level, both axial and polar Ricci carriers have
nonzero future-horizon-regular data and propagate on the Einstein
characteristics with admissible leading outer falloff.  The tested local
endpoint diagnostics therefore fail to select the Einstein image in either
parity.  The global populated Krein and connection theorems of this paper are
axial only.
\end{corollary}

\section{Logical map and evidence states}

The core structure is
\[
\begin{array}{ccccccccc}
0&\longrightarrow&M_{\rm RW}^{(2)}
&\longrightarrow&M_{\rm Bach}^{\rm axial}
&\overset{\delta\Ric}{\longrightarrow}&M_{\rm car}
&\longrightarrow&0\\[1mm]
&&\rotatebox{90}{$\subset$}&&&&
\rotatebox{90}{$\supset$}&&\\[-1mm]
&&F_1&\subset&F_2&\subset&F_3=M_{\rm axial},&&
\end{array}
\]
with successive factors \((2,2,1)\).  For real frequency, horizon data
\(h\in\C^3\) maps to
\[
\begin{array}{ccc}
&h&\\[-1mm]
T_- \swarrow && \searrow T_+\\[-1mm]
\mathscr I^-_{\rm in}&&\mathscr I^+_{\rm out},
\end{array}
\]
and the action current supplies
\[
H_{\mathcal H^+}+T_+^\dagger G_+T_+-T_-^\dagger G_-T_-=0.
\]

\begin{center}
\begin{tabularx}{0.98\textwidth}{@{}Y Y Y@{}}
\toprule
Statement & Evidence state & Boundary\\
\midrule
Ricci--Bach composition and exact sequence &
Exact symbolic theorem &
Schwarzschild, linearized, declared mode classes\\
RW/RW/spin-one filtration and partial jet &
Exact symbolic theorem with independent certificate &
Axial \(\ell=2\)\\
Incoming factor Gram and \(T_-\in GL(3,\C)\) &
Exact plus computer-assisted theorem &
Every real \(\omega>0\), \(M=1\) normalization\\
Outgoing population and pseudo-isometry &
Validated interval theorem &
\(I_*\); genericity off a locally finite set\\
Exact scalar threshold nonresonance &
Exact symbolic theorem &
No punctured zero-free scattering interval claimed\\
No growing separated mode &
Analytic filtered energy theorem &
Axial \(\ell=2\), declared complex boundary classes\\
Connection Smith valuations \((0,0,2)\) &
Validated projective Evans and local-unit theorem &
One enclosed simple damped QNM\\
Second-order Green-resolvent pole &
Open conditional corollary &
Requires analytic Fredholm realization\\
Time-domain boundedness and decay &
Open &
Requires threshold, limiting-absorption and resolvent bounds\\
\bottomrule
\end{tabularx}
\end{center}

\section{Conclusions}

The additional pure-Weyl carrier is not removed by the ordinary one-ended
Lorentzian Schwarzschild endpoint conditions tested here.  The result is
stronger than a local degree-of-freedom count:
\begin{itemize}
\item the carrier belongs to an exact non-split filtered module;
\item it enters the current derived from the Weyl action;
\item the complete incoming flux space is nondegenerate and indefinite;
\item every incoming direction is populated by a future-horizon-regular
solution at every positive real frequency;
\item outgoing population holds on a certified band and generically on the
positive axis;
\item the full filtered system has no growing separated mode;
\item one physical spin-two QNM is a defective connection-level double root.
\end{itemize}

Thus negative classical flux directions coexist with separated-mode
stability.  The remaining threat to time-domain control is quantitative
resolvent or pseudospectral growth, not an exponentially growing separated
frequency.

The next two calculations are sharply defined.  First, compute the explicit
typed \(T_+\) and reflection matrix over a substantial real-frequency band
using the correlated partial-jet transport.  Second, construct the analytic
Fredholm QNM realization needed to promote Theorem~\ref{thm:ep2} to a
second-order Green-resolvent pole.

\appendix

\section{Typed endpoint frames and explicit Grams}
\label{app:typed-grams}

The triangular factor frame and the diagonal signature frame serve different
purposes.  After normalizing the repeated-spin-two hyperbolic plane, write
\[
\mathcal B_0=(E,R,S),\qquad
J_0=
\begin{pmatrix}
0&1&0\\
1&0&0\\
0&0&-1
\end{pmatrix}.
\]
The corresponding signature frame is
\[
P=\frac{E+R}{\sqrt2},\qquad
N=\frac{E-R}{\sqrt2},\qquad
S=S,
\]
with
\[
J_\sigma=\operatorname{diag}(1,-1,-1),\qquad
C=
\begin{pmatrix}
1/\sqrt2&1/\sqrt2&0\\
1/\sqrt2&-1/\sqrt2&0\\
0&0&1
\end{pmatrix},
\quad C^\dagger J_0C=J_\sigma.
\]
If \(M_0\) is written in the factor frame, then
\[
M_\sigma=C^{-1}M_0C.
\]
For
\[
M_0=
\begin{pmatrix}
m&x&y\\0&m&z\\0&0&n
\end{pmatrix},
\]
one finds
\[
M_\sigma=
\begin{pmatrix}
m+x/2&-x/2&(y+z)/\sqrt2\\
x/2&m-x/2&(y-z)/\sqrt2\\
0&0&n
\end{pmatrix}.
\]
Thus triangularity is a statement about the filtration frame, not the
signature frame.  In particular a nonzero extension entry \(x\) produces
both upper and lower off-diagonal entries after signature normalization.

In the original incoming order
\((XI0,XI1,EI0)\), the normalized action Gram is
\begin{equation}
\label{eq:raw-gminus}
G_-^{\rm raw}=
\begin{pmatrix}
\dfrac{576}{5\omega}&\dfrac{96i}{5}&\dfrac{384\omega}{5}\\
-\dfrac{96i}{5}&-\dfrac{144\omega}{5}&-\dfrac{192i\omega^2}{5}\\
\dfrac{384\omega}{5}&\dfrac{192i\omega^2}{5}&0
\end{pmatrix}.
\end{equation}
The factor shear and orthogonal quotient lift reduce
\eqref{eq:raw-gminus} to \eqref{eq:factor-gram}.

For reference, in the original outgoing order
\((XI2,XI3,EI2)\), the normalized future Gram is
\begin{equation}
\resizebox{0.98\textwidth}{!}{$
G_+^{\rm raw}=
\begin{pmatrix}
\dfrac{384(-512\omega^6+640\omega^4-278\omega^2+39)}{5\omega}
&
\dfrac{12288i\omega^4-3072\omega^3-9984i\omega^2}{5}
+192\omega+384i
&
\dfrac{-1536i\omega^2+384\omega+768i}{5}
\\[1mm]
\dfrac{-12288i\omega^4-3072\omega^3+9984i\omega^2}{5}
+192\omega-384i
&
-\dfrac{768\omega^3}{5}+48\omega
&
\dfrac{96\omega}{5}
\\[1mm]
\dfrac{1536i\omega^2+384\omega-768i}{5}
&
\dfrac{96\omega}{5}
&
0
\end{pmatrix}.
$}
\end{equation}
The full physical Grams are \(\pi\alpha_{\rm W}G_\pm^{\rm raw}\).
On \([1/2,3/4]\), exact rational/Sturm bounds give
\[
\begin{array}{c|cc}
&\max\|G_\pm\|_F^2&\max\|G_\pm^{-1}\|_F^2\\ \hline
G_-&1429056/25&7025/65536\\
G_+&10389888/25&19825/65536 .
\end{array}
\]
These estimates give uniform equivalence of the finite-band coefficient
norms without inferring inverse bounds from determinants alone.

\subsection{The quotient metric is intrinsic}

Let \(K\) be the repeated-spin-two plane and \(M_1\) the spin-one quotient.
Because \(G_-|_K\) is nondegenerate,
\[
\pi_1:K^{\perp_{G_-}}\longrightarrow M_1
\]
is an isomorphism.  If \(s_\perp\) denotes its inverse, then
\[
g_1(q_1,q_2)=G_-(s_\perp q_1,s_\perp q_2)
\]
is independent of every arbitrary lift.  In a nonorthogonal adapted frame it
is the Schur complement
\[
g_1=G_{11}-G_{1K}G_{KK}^{-1}G_{K1}.
\]
The value \(g_1=-32\pi\alpha_{\rm W}/(15\omega)\) is therefore a
quotient-theoretic statement, not a sign attached to one chosen Witt vector.

\section{Krein scattering identities}
\label{app:krein-scattering}

In either typed frame define the Krein adjoint
\[
M^\sharp=J^{-1}M^\dagger J.
\]
After separate endpoint normalizations, the one-sided conservation law is
\[
R^\sharp R+A^\sharp A=I.
\]
The horizon defect operator and its Hermitian form are
\[
\mathfrak D_H=I-R^\sharp R=A^\sharp A,
\qquad
D_H=J-R^\dagger JR=A^\dagger JA.
\]
On the certified cell \(A=T_-^{-1}\) is invertible and the horizon Gram has
inertia \((1,2,0)\).  Hence
\[
\mathfrak D_H\in GL(3,\C),\qquad
1\notin\sigma(R^\sharp R),\qquad
\inertia D_H=(1,2,0).
\]
The defect is nondegenerate but indefinite.  It is not a positive absorption
operator, and the conservation theorem supplies no inequality of the form
``outgoing flux is at most incoming flux'' on the complete Weyl trace space.

In canonical normalized frames \(\det J=1\), so
\[
\det D_H=|\det A|^2=|\det T_-|^{-2}.
\]
Using
\[
\det T_-=A_{{\rm in},2}^2A_{{\rm in},1}
\]
and the scalar Wronskians,
\[
\det D_H=
|A_{{\rm in},2}|^{-4}|A_{{\rm in},1}|^{-2},
\qquad
0<\det D_H\le1.
\]
This positive signed three-volume is the surviving determinant-level remnant
of ordinary scalar absorption; the individual defect eigenvalues need not
lie in \([0,1]\).

An algebraic two-sided completion also follows.  The range of
\(\mathscr S=(R,A)^T\) is nondegenerate with inertia \((1,2)\) in a target of
inertia \((2,4)\).  Its orthogonal complement has inertia \((1,2)\).
Appending a \(J\oplus J\)-orthonormal basis of that complement gives an
element of \(U(2,4)\).  This is only an algebraic completion; it does not
construct the physical past-horizon channels.

\section{Projective Evans reduction and the certified disk}
\label{app:projective-evans}

For a scalar companion system
\[
Y_x=
\begin{pmatrix}A_{11}&A_{12}\\A_{21}&A_{22}\end{pmatrix}Y
\]
use the projective coordinate \(v=Y_2/Y_1\).  It satisfies
\[
v_x=A_{21}+(A_{22}-A_{11})v-A_{12}v^2.
\]
Let \(v_H\) and \(v_\infty\) be the horizon-regular and infinity-outgoing
lines transported to a common section.  Their mismatch
\[
\delta=v_\infty-v_H
\]
obeys the exact scalar equation
\[
\delta_x=
\bigl[A_{22}-A_{11}-A_{12}(v_\infty+v_H)\bigr]\delta.
\]
It follows that projective Evans functions at two matching sections differ
by an analytic nonzero factor.  Their zero divisors, logarithmic contour
moments, and local velocity ratios are match-section independent.

If \(\tau\) denotes the intrinsic partial-jet deformation, the tangent
variables satisfy
\[
(v_p)_x=F_vv_p+F_p,\qquad p\in\{\omega,\tau\},
\]
where \(F\) is the Riccati right-hand side.  At a simple zero,
\[
\kappa_n=
\frac{b(\omega_n)}{a'(\omega_n)}
=\frac{\delta_\tau(\omega_n)}{\delta_\omega(\omega_n)}.
\]
This removes every homogeneous amplitude from the Smith selector.

On a boundary contour avoiding zeros,
\[
\frac1{2\pi i}\oint\frac{\delta_\omega}{\delta}\,\dd\omega
\]
counts scalar QNMs, while
\[
\frac1{2\pi i}\oint\frac{\delta_\tau}{\delta}\,\dd\omega
\]
is the sum of their intrinsic velocity classes.  A projective chart change is
a M\"obius transformation.  It multiplies \(\delta\) by an analytic nonzero
unit whenever both endpoint lines remain in the chart, so the closed-contour
invariants are unchanged.

The certified disk in Theorem~\ref{thm:ep2} was closed in two stages.  First,
validated projective panels excluded zero on the boundary and supplied a
lifted argument change of \(2\pi\).  Second, a local interval-Newton rail
excluded zero from \(\delta_\omega\) and enclosed the unique simple root.
The intrinsic tangent and the independent spin-one mismatch were then
evaluated on the same disk.  This separation prevents a numerical root
approximation from being reused as its own simplicity or local-unit proof.

\section{Certificate authorities}

The principal machine authorities are:
\begin{itemize}
\item \cert{black_hole_programme/certificates/BH2A_AXIAL_OPERATOR.json};
\item \cert{black_hole_programme/phase3/axial_rw_lx_triangular_preflight/certificate.json};
\item \cert{black_hole_programme/phase3/axial_null_flux_gram/certificate.json};
\item \cert{black_hole_programme/phase3/axial_incoming_extended_domain_audit/certificate.json};
\item \cert{black_hole_programme/phase3/axial_boundary_devissage_no_growth/certificate.json};
\item \cert{black_hole_programme/phase3/axial_outgoing_population_cell_half_v1/certificate.json};
\item \cert{black_hole_programme/phase3/axial_global_finite_flux_channel_classification_v3/certificate.json};
\item \cert{black_hole_programme/phase4/axial_threshold_exact_structure_v1/certificate.json};
\item \cert{black_hole_programme/phase3/axial_qnm_projective_evans_contour_completion/full_contour_winding_v1/certificate.json};
\item \cert{black_hole_programme/phase3/axial_qnm_projective_evans_contour_completion/local_selector_v1/certificate.json};
\item \cert{black_hole_programme/phase3/axial_qnm_spin_one_local_unit_v1/certificate.json};
\item \cert{black_hole_programme/certificates/BH2B_POLAR_REACH.json}.
\end{itemize}
Each numerical theorem has a separately implemented verifier and mutation
tests.  Reproduction commands and content hashes are recorded in the generated
claim map and receipt accompanying this paper.

\section{What this paper does not establish}

It does not establish:
\begin{itemize}
\item explicit \(T_+\) entries or reflection amplitudes over a substantial
frequency band;
\item absence of isolated positive-real reflection zeros;
\item a punctured threshold interval of outgoing invertibility or the formal
low-frequency Jost and extension-ratio asymptotics in
Section~\ref{sec:threshold};
\item a complete asymptotically flat metric/tetrad phase space and charge
algebra;
\item a polar global connection theorem;
\item an analytic Fredholm realization of the damped QNM boundary problem;
\item a physical second-order Green-resolvent pole or generalized ringdown
term;
\item limiting absorption, uniform resolvent bounds, threshold control,
bounded time evolution, decay, or completeness;
\item a positive quantum metric, particles, a BRST physical Hilbert space,
CPT, or unitarity.
\end{itemize}

\begin{thebibliography}{99}

\bibitem{Maldacena}
J.~Maldacena,
\emph{Einstein gravity from conformal gravity},
arXiv:1105.5632 (2011).

\bibitem{ReggeWheeler}
T.~Regge and J.~A.~Wheeler,
\emph{Stability of a Schwarzschild singularity},
Phys.\ Rev.\ \textbf{108}, 1063 (1957).

\bibitem{LeeWald}
J.~Lee and R.~M.~Wald,
\emph{Local symmetries and constraints},
J.\ Math.\ Phys.\ \textbf{31}, 725 (1990).

\bibitem{IyerWald}
V.~Iyer and R.~M.~Wald,
\emph{Some properties of Noether charge and a proposal for dynamical black
hole entropy},
Phys.\ Rev.\ D \textbf{50}, 846 (1994).

\bibitem{Antoniou2025}
G.~Antoniou, L.~Gualtieri, and P.~Pani,
\emph{Gravitational quasinormal modes of black holes in quadratic gravity},
Phys.\ Rev.\ D \textbf{111}, 064059 (2025).

\bibitem{Knoska2025}
S.~Knoska, D.~Kofro\v n, and R.~\v Svarc,
\emph{Linear perturbations of Kerr black hole in quadratic gravity},
Phys.\ Rev.\ D \textbf{112}, 104058 (2025);
arXiv:2507.21923.

\bibitem{Konoplya2025}
R.~A.~Konoplya, A.~Spina, and A.~Zhidenko,
\emph{Time evolution of black hole perturbations in quadratic gravity},
Phys.\ Rev.\ D \textbf{112}, 024060 (2025);
arXiv:2505.01128.

\bibitem{Antoniou2026}
G.~Antoniou, L.~Gualtieri, and P.~Pani,
\emph{Nonperturbative suppression of beyond-general-relativity effects in
quadratic gravity},
Phys.\ Rev.\ D \textbf{113}, 124068 (2026).

\bibitem{Tajcovsky2026}
J.~V.~Taj\v covsk\'y,
\emph{Exact and Perturbative Solutions to Quadratic Gravity},
Master's thesis, Charles University (2026),
\url{https://hdl.handle.net/20.500.11956/209922}.

\end{thebibliography}

\end{document}
